\documentclass[a4paper, 12pt]{article}
\usepackage[margin=1in]{geometry}
\usepackage{graphicx}
\usepackage{amsfonts, amssymb, amsmath, amsthm}
\usepackage{tikz}
\usepackage{minted}

\DeclareMathOperator{\Ker}{Ker}
\DeclareMathOperator{\Image}{Im}

\renewcommand{\Im}{\Image}

\newcommand{\ceil}[1]{\lceil #1 \rceil}
\newcommand{\floor}[1]{\lfloor #1 \rfloor}
\newcommand{\lt}{\left}
\newcommand{\rt}{\right}
\newcommand{\C}{\mathbb{C}}
\newcommand{\R}{\mathbb{R}}
\newcommand{\Q}{\mathbb{Q}}
\newcommand{\Z}{\mathbb{Z}}
\newcommand{\N}{\mathbb{N}}
\newcommand{\normalsubgroup}{\lhd}
\newcommand{\cyclic}[1]{\langle #1 \rangle}
\newcommand{\epf}{\hfill \qed}
\newcommand{\chap}[1]{
\section*{#1}
\addcontentsline{toc}{section}{#1}
}
\newcommand{\ex}[1]{
\subsection*{#1}
\addcontentsline{toc}{subsection}{#1}
}

\begin{document}

\tableofcontents

\chap{Chapter 1 Groups and Homomorphisms}

\ex{Exercise 1}

Let $N$ be a subgroup of $G$ and $g \in G$. Then $g^{-1}Ng = \{g^{-1}ng \mid n \in N\} = \{g^{-1}gn \mid n \in N\} = \{n \mid n \in N\} = N$ because $G$ is abelian. So any subgroup of $G$ is normal. Since $G$ is simple, the only subgroups of $G$ are $\{1\}$ and $G$. Also, using the fact that $G$ is simple we can assume that $G \neq \{1\}$. Choose $h \in G \setminus \{1\}$. Since $\cyclic{h}$ is a subgroup of $G$, $\cyclic{h} = G$ or $\cyclic{h} = \{1\}$. If $\cyclic{h} = \{1\}$, then $h=1$, a contradiction. Hence, $\cyclic{h} = G$, so $G$ is a cyclic group. Let $|G|=n$. Assume for contradiction that $n$ is not prime, so $n=ab$ for some $1<a,b<n$. Since $G$ is cyclic and $a$ is a factor of $n$, there exists $k \in G$ of order $a$. Since $|k|>1$, $k \neq 1$. By the above argument, we have $\cyclic{k}=G$. Therefore, $a = |k| = |\cyclic{k}| = |G| = n$, a contradiction. Hence, $G$ is cyclic of prime order. \epf

\ex{Exercise 2}

Assume $H \neq \{1\}$. Let $g_1, g_2 \in G$. Suppose $g_1\vartheta = g_2\vartheta$. Then $(g_1\vartheta)(g_2\vartheta)^{-1}=1$ and thus $(g_1g_2^{-1})\vartheta=1$. Hence $g_1g_2^{-1} \in \Ker(\vartheta)$. Since $\Ker(\vartheta) \normalsubgroup G$, we have $\Ker(\vartheta) = \{1\}$ or $\Ker(\vartheta) = G$. If $\Ker(\vartheta) = G$, then $g\vartheta=1$ for all $g \in G$. Since $\vartheta$ is surjective, $H = \Im(\vartheta) = \{g\vartheta \mid g \in G\} = \{1\}$, a contradiction. Therefore, $\Ker(\vartheta)=\{1\}$. Hence, we have $g_1g_2^{-1}=1$, which implies $g_1=g_2$. So $\vartheta$ is injective and thus is bijective. \epf

\chap{Chapter 2 Vector Spaces and Linear Transformations}

\ex{Exercise 3}

Suppose $V = U \oplus W$. Clearly, $V=U+W$. Let $v \in U \cap W \subseteq V$. Then there exist unique $u \in U$ and $w \in W$ such that $v=u+w$. So $u+w \in U$ and $u+w \in W$. By the uniqueness of $u$ and $w$, we have $u+w=u$ and $u+w=w$, so $u=w=0$, and thus $v=0$. Hence, $U \cap W = \{0\}$.

Now suppose $V=U+W$ and $U \cap W = \{0\}$. We want to prove that for any $v \in V$ there exist unique $u \in U$ and $w \in W$ such that $v=u+w$. Let $v \in V$. Suppose $u_1,u_2 \in U$ and $w_1,w_2 \in W$ satisfy $v=u_1+w_1=u_2+w_2$. Then $u_1-u_2=w_2-w_1$. Since $U$ and $W$ are subspaces, they are closed under vector addition and additive inverses, so $u_1-u_2 \in U$ and $w_2-w_1 \in W$, so $u_1-u_2=w_2-w_1 \in U \cap W$. Hence, $u_1-u_2=w_2-w_1=0$, and thus $u_1=u_2$ and $w_1=w_2$, proving the uniqueness of $u$ and $w$.

Therefore, $V = U \oplus W$ if and only if $V=U+W$ and $U \cap W = \{0\}$. \epf

\ex{Exercise 4}

Suppose $V = U \oplus W$. Let $v \in V$. Then there exist unique $u \in U$ and $w \in W$ such that $v=u+w$. Since $u_1,\ldots,u_r$ is a basis of $U$, there exist unique $\lambda_1,\ldots,\lambda_r \in F$ such that $u = \lambda_1u_1 + \ldots + \lambda_ru_r$. Since $w_1,\ldots,w_s$ is a basis of $W$, there exist unique $\mu_1,\ldots,\mu_s \in F$ such that $w = \mu_1w_1 + \ldots + \mu_sw_s$. Therefore, there exist unique $\lambda_1,\ldots,\lambda_r,\mu_1,\ldots,\mu_s \in F$ such that $v=\lambda_1u_1+\ldots+\lambda_ru_r+\mu_1w_1+\ldots+\mu_sw_s$. Hence, $u_1,\ldots,u_r,w_1,\ldots,w_s$ is a basis of $V$.

Now suppose $u_1,\ldots,u_r,w_1,\ldots,w_s$ is a basis of $V$. Let $v \in V$. Then there exist unique $\lambda_1,\ldots,\lambda_r,\mu_1,\ldots,\mu_s \in F$ such that $v = \lambda_1u_1 + \ldots + \lambda_ru_r + \mu_1w_1 + \ldots + \mu_sw_s$. Since $u_1,\ldots,u_r$ is a basis of $U$, there exists a unique $u \in U$ such that $u=\lambda_1u_1 + \ldots + \lambda_ru_r$. Since $w_1,\ldots,w_s$ is a basis of $W$, there exists a unique $w \in W$ such that $w = \mu_1w_1 + \ldots + \mu_sw_s$. Hence, there exist unique $u \in U$ and $w \in W$ such that $v=u+w$. So $V = U \oplus W$.

Therefore, $V = U \oplus W$ if and only if $u_1,\ldots,u_r,w_1,\ldots,w_s$ is a basis of $V$. \epf



\end{document}
